% Figure: Schematic of a single cell on a substrate, with focal
% adhesions, actin stress fibres, and the force balance through the
% integrin-ECM coupling.
\begin{figure}[htbp]
\centering
\begin{tikzpicture}[
  axis/.style={-{Stealth}, thick, draw=engblack},
  ecm/.style={engred!30, draw=engred!70, thick},
  cell/.style={engblue!15, draw=engblue!70, very thick},
  stressfibre/.style={engred, thick},
  fa/.style={engblack, very thick, fill=engblack},
  annot/.style={font=\scriptsize, align=left},
]

% Substrate (ECM)
\draw[ecm, fill=engred!10] (-0.5, -0.5) rectangle (10, 0.5);
\node[annot, anchor=north] at (5.0, -0.6) {Substrate (extracellular matrix, $E_{\text{sub}}$)};

% Cell body: flattened ellipse
\draw[cell] (5, 0.5) ellipse (3.5 and 1.6);

% Nucleus
\draw[engblue!40, draw=engblue!70, thick, fill=engblue!20] (5, 1.0) ellipse (1.0 and 0.5);
\node[font=\scriptsize] at (5, 1.0) {N};

% Focal adhesions (small black rectangles at the basal surface)
\foreach \x in {2.5, 3.4, 4.2, 5.0, 5.8, 6.6, 7.5} {
  \fill[fa] (\x - 0.12, 0.4) rectangle (\x + 0.12, 0.55);
}

% Stress fibres (lines connecting focal adhesions across cell)
\draw[stressfibre] (2.5, 0.5) -- (7.5, 0.5);
\draw[stressfibre] (3.4, 0.5) -- (6.6, 0.5);
\draw[stressfibre] (4.2, 0.5) -- (5.8, 0.5);
\draw[stressfibre] (2.5, 0.5) .. controls (3.5, 1.4) and (6.5, 1.4) .. (7.5, 0.5);

% Traction force arrows: cells pull the substrate inward
\foreach \x/\d in {2.5/-0.5, 3.4/-0.3, 4.2/-0.1, 5.8/0.1, 6.6/0.3, 7.5/0.5} {
  \draw[-{Stealth}, very thick, engblue] (\x, 0.2) -- (\x + \d, 0.2);
}

% Labels
\node[annot, anchor=west] at (8.3, 0.5) {focal\\adhesion};
\draw[-{Stealth}, thin] (8.2, 0.55) -- (7.6, 0.55);

\node[annot, anchor=west] at (8.3, 1.5) {actin stress\\fibre};
\draw[-{Stealth}, thin] (8.2, 1.4) -- (6.5, 0.9);

\node[annot, anchor=east] at (1.5, 1.4) {cell body\\(cortex + cytoplasm)};
\draw[-{Stealth}, thin] (1.6, 1.4) -- (2.5, 1.5);

\node[annot, anchor=west] at (5.0, -0.2) {traction force $\sim 1\,\text{nN}$ per FA};

% Top: stress-fibre tension annotation
\node[annot, engred] at (5.0, 2.4)
  {Stress-fibre tension $T \sim 10$ to $100\,\text{nN}$};

% Sum of traction at focal adhesions
\node[annot, anchor=west] at (-0.5, 3.5)
  {$\displaystyle \sum_i \vec{T}_i = 0$ (cell-substrate force balance)\\
   $\displaystyle \sum_i \vec{T}_i \times \vec{r}_i = 0$ (moment balance)};

\end{tikzpicture}
\caption{Force balance on an adherent cell. The cell exerts traction
forces (blue arrows) on the substrate at focal adhesion plaques,
through transmembrane integrin receptors that link the substrate's
extracellular matrix to the cell's actin cytoskeleton. Internal actin
stress fibres span between focal adhesions and carry tension
$T \sim 10$ to $100\,\text{nN}$ generated by myosin-II motors. The
tractions must sum to zero (force balance) and their moments must
sum to zero (moment balance), because the cell is in quasi-static
equilibrium. Traction-force microscopy measures the tractions by
tracking displacements of beads embedded in an elastic substrate.}
\label{fig:vol06ch06:cell-force}
\end{figure}
